% STATUS: rewritten 2026-11-05; every claim is cited.
\section{BCK, BCI and MV-algebras}\label{sec:other-systems:bci-mv}

Here a BCI or BCK algebra keeps only an operation.
This operation behaves like the difference
$x - y$ of a group or the implication $y \to x$ of a logic
\cite[p.~167]{JunXin2004}. Jun and Xin
introduced derivations on BCI-algebras in 2004
\cite[p.~167]{JunXin2004}.

The derivations below act on the subtraction and implication operations of a BCI or MV algebra, which is what the product rule becomes in that setting.

\begin{definition}[{\cite[p.~168]{JunXin2004}}]\label{def:other:bci}
A nonempty set $X$ with a constant $0$ and a binary operation, written by
juxtaposition, is a \emph{BCI-algebra} if, for all $x, y, z \in X$,
\textup{(I)} $((xy)(xz))(zy) = 0$,
\textup{(II)} $(x(xy))y = 0$,
\textup{(III)} $xx = 0$, and
\textup{(IV)} $xy = 0$ and $yx = 0$ imply $x = y$.
A BCI-algebra $X$ with $0 \le x$ for all $x \in X$, where $x \le y$ means
$xy = 0$, is a \emph{BCK-algebra} \cite[p.~168]{JunXin2004}.
\end{definition}

The next definition states the derivations on a BCI-algebra.

\begin{definition}[{\cite[Definition~3.5, p.~170]{JunXin2004}}]\label{def:other:lr-derivation}
Let $X$ be a BCI-algebra and write $x \wedge y = y(yx)$. A self-map $d$
of $X$ is a \emph{left-right derivation} (an \emph{$(l, r)$-derivation})
of $X$ if
\[
d(xy) = d(x)y \wedge x\,d(y) \qquad (x, y \in X),
\]
and a \emph{right-left derivation} (an \emph{$(r, l)$-derivation}) if
$d(xy) = x\,d(y) \wedge d(x)y$ for all $x, y \in X$.
\end{definition}

A derivation $d$ on a BCI-algebra is \emph{regular} if $d(0) = 0$
\cite[Definition~3.12, p.~173]{JunXin2004}. Regular derivations detect
the p-semisimple BCI-algebras (Theorem~\ref{thm:other:jun-xin-317})
\cite[Theorem~3.17, p.~174]{JunXin2004}.\begin{theorem}[{\cite[Theorem~3.17, p.~174]{JunXin2004}}]\label{thm:other:jun-xin-317}
Let $X$ be a BCI-algebra. Then the following are equivalent:
\textup{(i)} $X$ is p-semisimple; \textup{(ii)} if $X$ has a regular
derivation, then it is trivial; \textup{(iii)} $d^{-1}(0) = \{0\}$ for
every regular derivation $d$ of $X$.
\end{theorem}

The next theorem gives an ideal-theoretic test for regularity (Theorem~\ref{thm:other:jun-xin-322}).

\begin{theorem}[{\cite[Theorem~3.22, p.~176]{JunXin2004}}]\label{thm:other:jun-xin-322}
Let $d$ be a derivation of a BCI-algebra $X$. Then $d$ is regular if and
only if every ideal of $X$ is $d$-invariant.
\end{theorem}

Ciungu widened the class to the pseudo-BCK algebras, where the single
arrow splits into two
\cite[Definition~2.1, p.~177]{Ciungu2021}.

\begin{definition}[{\cite[Definition~2.1, p.~177]{Ciungu2021}}]\label{def:other:pseudo-bck}
A \emph{pseudo-BCK algebra} is a structure $(A, \le, \to, \rightsquigarrow, 1)$
where $\le$ is a binary relation on $A$, $\to$ and $\rightsquigarrow$ are
binary operations on $A$ and $1 \in A$, satisfying, for all $x, y, z \in A$:
$x \to y \le (y \to z) \rightsquigarrow (x \to z)$ and
$x \rightsquigarrow y \le (y \rightsquigarrow z) \to (x \rightsquigarrow z)$;
$x \le (x \to y) \rightsquigarrow y$ and
$x \le (x \rightsquigarrow y) \to y$; $x \le x$; $x \le 1$; if $x \le y$
and $y \le x$ then $x = y$; and $x \le y$ if and only if $x \to y = 1$ if
and only if $x \rightsquigarrow y = 1$.
\end{definition}

The next definition states the implicative derivations on this algebra.

\begin{definition}[{\cite[pp.~178--179]{Ciungu2021}}]\label{def:other:implicative-derivation}
Let $(A, \to, \rightsquigarrow, 1)$ be a pseudo-BCK algebra. A map
$d \colon A \to A$ is a \emph{type I implicative derivation} on $A$ if,
for all $x, y \in A$,
\[
d(x \to y) = (x \to d(y)) \vee_1 (d(x) \to y),
\qquad
d(x \rightsquigarrow y) = (x \rightsquigarrow d(y)) \vee_2 (d(x) \rightsquigarrow y),
\]
and a \emph{type II implicative derivation} if the two identities hold
with the summands on the right interchanged, where $\vee_1$ and
$\vee_2$ are the two joins of the algebra.
\end{definition}

An MV-algebra is the model of Łukasiewicz logic that carries both an
addition and a logical not \cite[p.~261]{WangSheQian2019}.

\begin{definition}[{\cite[p.~261]{WangSheQian2019}}]\label{def:other:mv-algebra}
An \emph{MV-algebra} is an algebra $L = (L, \oplus, {}^{*}, 0)$
satisfying, for all $x, y \in L$:
\textup{(MV1)} $(L, \oplus, 0)$ is a commutative monoid;
\textup{(MV2)} $(x^{*})^{*} = x$;
\textup{(MV3)} $0^{*} \oplus x = 0^{*}$;
\textup{(MV4)} $(x^{*} \oplus y)^{*} \oplus y = (y^{*} \oplus x)^{*} \oplus x$.
Write $1 = 0^{*}$ and $x \to y = x^{*} \oplus y$.
\end{definition}

The next definition states the derivations that respect the implication.

\begin{definition}[{\cite[Definition~3.1, p.~262]{WangSheQian2019}}]\label{def:other:mv-implicative}
Let $L$ be an MV-algebra. A mapping $g \colon L \to L$ is an
\emph{implicative derivation} on $L$ if, for all $x, y \in L$,
\[
g(x \to y) = (g(x) \to y) \vee (x \to g(y)).
\]
The set of all implicative derivations of $L$ is denoted $I(L)$.
\end{definition}

On a Boolean algebra, the set $I(L)$ is itself a Boolean algebra, and the
original algebra appears as its algebra of implicative derivations (Theorem~\ref{thm:other:wang-310})
\cite[Theorem~3.10, p.~267]{WangSheQian2019}.\begin{theorem}[{\cite[Theorem~3.10, p.~267]{WangSheQian2019}}]\label{thm:other:wang-310}
Let $(L, \wedge, \vee, {}^{*}, 0, 1)$ be a Boolean algebra. Then
$(I(L), \cap, \cup, \neg, g_{1}, 1_{g})$, with $g_{1}(x) = x$ and
$1_{g}(x) = 1$, is a Boolean algebra.
\end{theorem}

The next theorem recovers the original algebra from its derivations (Theorem~\ref{thm:other:wang-311}).

\begin{theorem}[{\cite[Theorem~3.11, p.~268]{WangSheQian2019}}]\label{thm:other:wang-311}
Let $(L, \wedge, \vee, {}^{*}, 0, 1)$ be a Boolean algebra. Then $L$ is
isomorphic to the Boolean algebra $(I(L), \cap, \cup, \neg, g_{1}, 1_{g})$
of all implicative derivations on $L$, by the map $a \mapsto \varphi(a)$,
where $\varphi(a)(x) = a \vee x$.
\end{theorem}

The next section takes up the last family of the chapter, the subtraction
algebras, where the meet rule takes the place of the product rule once
more \cite[Definition~3.1, p.~511]{OzturkCeven2009}.
